% Material richtig einlegen: Lage von Materialrändern und Andruckrollen.
% Eigene schematische Zeichnung (TikZ), keine Vorlage aus der Roland-Anleitung.
% Lizenz wie das Dokument: CC BY-SA 3.0
\begin{tikzpicture}[x=1cm, y=1cm, font=\sffamily\footnotesize,
	teil/.style={draw=black!75, line width=0.5pt},
	markierung/.style={fill=black!45},
	beschriftung/.style={align=left, inner sep=1pt, text width=4.2cm},
	zeiger/.style={draw=black!80, line width=0.4pt, -{Circle[length=2pt]}},
	]
	% Balken mit Kontrollmarkierungen
	\draw[teil, fill=black!20] (0,3.0) rectangle (12,3.6);
	\fill[markierung] (0.2,3.42) rectangle (2.0,3.56);            % große Markierung
	\foreach \x in {4.3,6.5,8.7,10.5}
		\fill[markierung] (\x,3.42) rectangle (\x+0.7,3.56);      % kleine Markierungen
	% Auflage und Schneidekante
	\draw[teil, fill=white] (0,0.6) rectangle (12,3.0);
	\fill[cyan!35] (0,1.3) rectangle (12,1.5);
	% Material: linker Rand in der großen Markierung, rechter Rand am rechten Ende einer kleinen
	\fill[blue!12, opacity=0.85] (0.6,0.25) rectangle (9.4,3.0);
	\draw[blue!60!black, line width=0.7pt] (0.6,0.25) -- (0.6,3.0) (9.4,0.25) -- (9.4,3.0)
		(0.6,0.25) -- (9.4,0.25);
	\node[blue!50!black, font=\sffamily\small] at (5.0,0.85) {Material};
	% Andruckrollen: links in der großen Markierung, rechts in einer kleinen, ganz auf dem Material
	\foreach \x in {1.3,9.05}
		\draw[teil, fill=black!55, rounded corners=0.6mm] (\x-0.25,2.75) rectangle (\x+0.25,3.25);
	% Hilfslinien vorn
	\draw[black!70, line width=0.8pt, dash pattern=on 2mm off 1mm] (0.2,0.05) -- (0.2,0.6);
	\draw[black!70, line width=0.8pt, dash pattern=on 2mm off 1mm] (11.8,0.05) -- (11.8,0.6);
	% Maßhilfe rechter Rand
	\draw[black!60, densely dotted] (9.4,3.6) -- (9.4,4.1);

	\node[beschriftung, anchor=south west] (gross) at (-0.6,4.3)
		{\textbf{Große Markierung}\\ linker Materialrand und linke Andruckrolle liegen darunter};
	\draw[zeiger] (gross.south west) ++(0.6,0) -- (1.0,3.5);
	\node[beschriftung, anchor=south west] (klein) at (4.4,4.3)
		{\textbf{Vier kleine Markierungen}\\ rechte Andruckrolle unter eine davon, Abstand der Rollen möglichst groß};
	\draw[zeiger] (klein.south west) ++(0.3,0) -- (4.65,3.5);
	\node[beschriftung, anchor=south west, text width=3.6cm] (rand) at (9.0,4.3)
		{\textbf{Rechter Materialrand}\\ genau unter einer Markierung, Rolle ganz auf dem Material};
	\draw[zeiger] (rand.south west) ++(0.3,0) -- (9.4,3.9);
	\node[beschriftung, anchor=north west] (hilf) at (-0.6,-0.5)
		{\textbf{Hilfslinien}\\ Material gerade einlegen, Vorderkante bis hierhin};
	\draw[zeiger] (hilf.north west) ++(0.6,0) -- (0.2,0.15);
	\node[beschriftung, anchor=north west, text width=3.2cm] (kante) at (9.6,-0.5)
		{\textbf{Schneidekante}\\ hellblauer Gummistreifen};
	\draw[zeiger] (kante.north west) ++(0.3,0) -- (10.6,1.4);
\end{tikzpicture}
